早泄并非单一疾病，而是一组表现相似、病因各异的综合征。目前国际公认的分类以 Waldinger 提出的四型框架最为实用，国际性医学会（ISSM）与欧洲泌尿外科学会（EAU）的指南均予采用：

\begin{itemize}
	\item \textbf{终身性（原发性）早泄}：自首次性生活起即持续存在，约 90\% 的射精发生在插入后 1 分钟内，几乎每次、几乎每种性情境都会发生；多与神经生物学因素（5-羟色胺能抑制不足、射精阈值先天偏低）主导有关。
	\item \textbf{获得性（继发性）早泄}：曾经具有正常的射精控制力，之后逐渐出现并加重；射精潜伏期常短于 3 分钟，或较本人既往水平明显缩短；常继发于勃起功能障碍、慢性前列腺炎、甲状腺功能亢进、焦虑抑郁或关系变故等，须查找并治疗原发病。
	\item \textbf{自然变异性早泄}：射精过快呈\textbf{偶发、情境性}，并非持续存在，控制力时好时坏，本人苦恼程度通常较轻。它属于正常范围内的波动，\textbf{一般无需药物治疗}。
	\item \textbf{早泄样射精功能障碍}：患者主观认定"自己早泄"，但射精潜伏期其实在正常范围（多在 3 分钟以上），控制感也不差；其困扰往往源于对"理想时长"的过高期待，或勃起维持困难时的"抢时间"心理，处理重点是心理教育与认知调整，而非一味加药。
\end{itemize}

\begin{tcolorbox}[colback=pink!5!white,colframe=red!50!black,title=贴心小叮咛]
后两型（自然变异性、早泄样功能障碍）最容易被误诊误治：前者本不必吃药，后者的症结不在"时间"而在"期待与信心"。先分清这四型，是避免过度医疗的第一步。
\end{tcolorbox}

\noindent\textbf{按严重程度划分}（便于沟通，非诊断金标准）：轻度——射精潜伏期约 1--3 分钟，能射但控制困难、偶有困扰；中度——约 30 秒至 1 分钟，控制力差、常伴明显苦恼；重度——多在 30 秒内甚至插入前即射精，几乎无法控制。

\noindent\textbf{关于"多长时间算早泄"}：终身性早泄的射精潜伏期通常短于 1 分钟，获得性早泄通常短于 3 分钟或较既往明显缩短。但必须强调，时间只是参考，"无法自主延迟 + 造成本人或伴侣的困扰、焦虑或回避"才是诊断的核心。
